\chapter{INTRODUCTION}
\label{ch:introduction}

Large Language Models (LLMs) have demonstrated unprecedented capabilities in natural language understanding, leading to their widespread deployment through Machine Learning as a Service (MLaaS) platforms. However, sending sensitive user data---such as medical records, legal documents, or proprietary corporate logic---to remote third-party servers raises severe privacy concerns. To mitigate these risks, advanced cryptographic techniques such as Fully Homomorphic Encryption (FHE) have emerged as a highly promising solution. FHE allows complex mathematical operations to be evaluated directly on ciphertexts, ensuring that the cloud server processes the prompt without ever accessing the plaintext data.

Adapting modern neural networks to operate within an encrypted domain is a highly non-trivial task. Standard LLM architectures were designed for plaintext hardware acceleration, not the strict mathematical constraints of cryptographic rings. Consequently, ensuring data privacy currently comes at the cost of massive computational overhead and algorithmic restructuring.

\section{Motivation}
Despite its theoretical elegance, integrating FHE with massive Transformer architectures introduces staggering memory requirements and latency bottlenecks. FHE schemes natively support only addition and multiplication, meaning the exact non-linear activation functions used in LLMs are fundamentally incompatible with the encrypted domain. The three primary non-linearities---Softmax in self-attention, GELU in feed-forward networks, and the inverse square root in LayerNorm---must each be replaced by low-degree polynomial approximations that consume a bounded number of multiplicative levels.

Existing approaches to this challenge---such as THE-X~\cite{chen2022thex}, Iron~\cite{hao2022iron}, MPCFormer~\cite{li2023mpcformer}, and BOLT~\cite{pang2024bolt}---replace only two of the three non-linear operations (typically GELU and Softmax), while either leaving LayerNorm intact, delegating it to client-side computation, or replacing it with a simple affine transformation that discards the normalisation effect entirely. This leaves a critical gap: LayerNorm cannot be evaluated homomorphically, preventing fully non-interactive encrypted inference.

Table~\ref{tab:plaintext_vs_fhe} illustrates the critical trade-offs between standard inference and FHE-based secure inference, highlighting the motivation for optimising these cryptographic systems.

\begin{table}[h]
\centering
\caption{Comparison of plaintext and FHE-based LLM inference paradigms.}
\label{tab:plaintext_vs_fhe}
\begin{tabular}{|l|c|c|}
\hline
\textbf{Metric} & \textbf{Plaintext Inference} & \textbf{FHE Inference (CKKS)} \\ \hline
Data Privacy & None (Server sees all) & Absolute (Server sees ciphertexts) \\ \hline
Activation Functions & Exact (Softmax, GELU, LN) & Approximate (Polynomials) \\ \hline
Inference Latency & Milliseconds per token & Seconds--minutes per token \\ \hline
Multiplicative Depth & Unlimited & Fixed budget ($D$ levels) \\ \hline
\end{tabular}
\end{table}

\subsection{Problem Statement}
The primary challenge addressed in this thesis is the complete replacement of all non-polynomial operations in the Transformer architecture with learnable polynomial approximations that can be evaluated under CKKS encryption, while maintaining competitive accuracy with the original model. The core research question is:

\begin{quote}
\emph{How can all three non-linear operations in Transformers---GELU, Softmax, and LayerNorm---be replaced by learnable polynomial approximations using a progressive training strategy, such that the resulting fully-polynomial model achieves accuracy competitive with the original while enabling non-interactive encrypted inference?}
\end{quote}


\section{Contributions}

This thesis proposes the \textbf{Learnable Polynomial Approximation Network (LPAN)} framework and makes the following five contributions:

\begin{enumerate}
    \item \textbf{LPAN: A Three-Stage Progressive Polynomial Replacement Framework.} A complete pipeline that replaces all three non-linear operations in BERT with learnable Chebyshev polynomials through a staged process: Stage~1 replaces GELU via cross-entropy fine-tuning, Stage~2 progressively replaces Softmax layer-by-layer via attention-level knowledge distillation, and Stage~3 progressively replaces LayerNorm layer-by-layer via the same distillation strategy with co-adaptation.

    \item \textbf{Distribution-Weighted Minimax Polynomial Initialisation.} A novel approximation method that incorporates the empirical activation distribution $\rho(x)$ into the minimax objective, concentrating approximation effort in high-density regions. This provides superior polynomial coefficient initialisation, achieving $3.4\times$ lower weighted error than standard Chebyshev approximation on GELU at degree~8. These initialised coefficients are subsequently refined through backpropagation during training.

    \item \textbf{Progressive Layer-by-Layer Replacement with Attention-Level Knowledge Distillation.} A training strategy for deep models that replaces operations one layer at a time, using a composite loss $\mathcal{L} = \alpha \mathcal{L}_{\mathrm{CE}} + \beta \sum \mathrm{KL}(A_{\mathrm{teacher}} \| A_{\mathrm{student}}) + \gamma \sum \mathrm{MSE}(H_{\mathrm{teacher}}, H_{\mathrm{student}})$ that preserves both attention patterns and hidden-state representations throughout the replacement process. Co-adaptation unfreezes all previously replaced layers during each sub-step, allowing joint refinement.

    \item \textbf{Per-Head Polynomial Softmax.} An architecture where each attention head maintains independent learnable Chebyshev coefficients for the exponential approximation, enabling head specialisation. Unlike shared-polynomial approaches, per-head coefficients capture the distinct attention patterns learned by different heads.

    \item \textbf{Complete Polynomial LayerNorm Replacement.} The first approach to replace LayerNorm's inverse square root $1/\sqrt{\sigma^2 + \epsilon}$ with a learnable polynomial approximation across all layers of a deep Transformer, maintaining competitive accuracy. Prior methods (THE-X, Iron, BOLT, MPCFormer) either skip LayerNorm replacement, delegate it to client-side computation, or replace it with a simple affine transformation.
\end{enumerate}

All contributions are validated through a complete end-to-end pipeline on four BERT model variants (Tiny, Mini, Small, and Base) on the SST-2 sentiment classification benchmark. The LPAN framework achieves:
\begin{itemize}
    \item \textbf{BERT-Tiny:} 83.26\% $\to$ 83.14\% (\textbf{0.11\% drop})
    \item \textbf{BERT-Mini:} 87.16\% $\to$ 86.81\% (\textbf{0.34\% drop})
    \item \textbf{BERT-Small:} 87.73\% $\to$ 88.53\% (\textbf{0.80\% gain})
    \item \textbf{BERT-Base:} 92.20\% $\to$ 91.86\% (\textbf{0.34\% drop})
\end{itemize}
These results surpass all existing methods that replace only two of three non-linear operations, while our approach replaces \emph{all three}---including the previously untouched LayerNorm.


\section{Thesis Organisation}

The remainder of this thesis is organised as follows:

\begin{itemize}
    \item \textbf{Chapter~\ref{ch:background}} provides background on the CKKS homomorphic encryption scheme, the non-linearity bottleneck in Transformers, and existing secure inference frameworks (CryptoNets, Iron, THE-X, MPCFormer, BOLT).

    \item \textbf{Chapter~\ref{ch:methodology}} presents the proposed LPAN methodology, including distribution-weighted minimax polynomial initialisation, the three-stage progressive replacement pipeline, attention-level knowledge distillation, per-head polynomial Softmax, polynomial LayerNorm, and error propagation analysis.

    \item \textbf{Chapter~\ref{ch:results}} reports experimental results across four BERT model variants, including activation profiling, approximation quality, progressive training analysis for BERT-Base, comparison with existing frameworks, and CKKS encrypted inference benchmarks.

    \item \textbf{Chapter~\ref{ch:conclusion}} summarises the contributions, discusses limitations, and outlines future research directions including GPU-accelerated CKKS inference and multi-task evaluation.
\end{itemize}
